% ==== AUTO-GENERADA — no editar a mano (regenerar con generate_thesis_tables.py) ====
\begin{table}[htbp]
\centering
\caption[Conteo de ejecuciones por modelo]{Conteo de ejecuciones del pipeline por modelo en la campaña experimental (fuente: ResultIndex, elaboración propia). Una \emph{ejecución} es una corrida completa de las Fases 1--3 para una (configuración, semilla).}
\label{tab:auto_campaign}
\begin{tabular}{lcc}
\toprule
Modelo & Ejecuciones & Tasa aprob. \\
\midrule
Heisenberg XXZ & 35 & 0/35 (0,0\%) \\
Heisenberg transversal & 40 & 0/40 (0,0\%) \\
Cadena de Kitaev & 7 & 0/7 (0,0\%) \\
TFIM & 490 & 241/490 (49,2\%) \\
TFIM por enlace & 511 & 464/511 (90,8\%) \\
TFIM frustrado ($J_1$-$J_2$) & 18 & 6/18 (33,3\%) \\
TFIM longitudinal & 185 & 55/185 (29,7\%) \\
\textbf{Total} & \textbf{1286} & \textbf{766/1286 (59,6\%)} \\
\bottomrule
\end{tabular}
\\[2pt]
\footnotesize Heisenberg: 35 XXZ + 40 transversal = 75 ejecuciones. La columna \emph{Tasa aprob.} agrega \emph{todas} las ejecuciones del modelo, incluidas las de fuera del régimen operativo válido, las cinco topologías y las cuatro profundidades; por tanto \emph{no} es comparable con las tasas por configuración de las Tablas~\ref{tab:cross_topo}, \ref{tab:cross_topo_depth} y \ref{tab:scaling} (restringidas al régimen válido y a la mejor profundidad). Es un indicador de cobertura de la campaña, no del rendimiento del pipeline en su régimen de operación.
\end{table}
